\chapter{S3：动态倾角与抗干扰验证}
\section{阶段定位与设备边界}
本阶段验证\textbf{动态倾角仪}在动态运动和干扰下的真实输出性能。对需纯线性加速度或振动激励的条目，目前双轴转台不足以单独提供激励，必须另配设备。动态角度参考、同步和采样延迟的不确定度需先评估。
\section{S3 测试项目总表}
\small
\begin{longtable}{C{18mm}p{46mm}p{61mm}p{22mm}}
\toprule 编号 & 测试项 & 测试目标／方法提示 & 级别／条件\\\midrule\endfirsthead
\toprule 编号 & 测试项 & 测试目标／方法提示 & 级别／条件\\\midrule\endhead
S3-00 & 动态参考/同步能力确认 & 动态角、速度、采集时基计量与转台限制 & 前置/必做\\
S3-01 & 正弦摆动幅值误差 & 多角度幅值与角速度上限联合设计 & P0\\
S3-02 & 动态跟踪精度 & 整段 RMSE、P95、最大瞬态误差 & P0\\
S3-03 & 幅频与相频特性 & 扫频、增益、相位曲线与带宽 & P0\\
S3-04 & 动态时延和时间戳可靠性 & 端到端延迟、数据时间对齐及不确定度 & P0\\
S3-05 & 阶跃/快速变位响应 & 超调、调整时间、滤波恢复 & P1\\
S3-06 & 多速率/变速/反向跟踪 & 匀速、加减速、方向反转残差 & P1\\
S3-07 & 双轴组合运动响应 & 按真实 U-T 运动学构造复合参考 & P1/有条件\\
S3-08 & 纯线性加速度干扰 & 实际姿态不变的水平/竖直加速度输入 & P0/需额外设备\\
S3-09 & 持续线性加速度和恢复 & 长期外加速度误判与停止后恢复 & P1/需额外设备\\
S3-10 & 旋转偏心/杠杆臂效应 & 传感器偏离旋转中心造成的加速度 & P1\\
S3-11 & 正弦/随机振动干扰 & 频谱、RMS 角度误差、振动整流 & P1/需振动台\\
S3-12 & 冲击与短时饱和恢复 & 冲击后的零点跳变、输出恢复 & P1/需相应设备\\
S3-13 & 丢帧/采样率变化的鲁棒性 & 速率变化与采集异常下姿态解算性能 & P1\\
S3-14 & 持续动态运行与漂移 & 重复周期运动、热平衡及系统偏差增长 & P1\\
S3-15 & 动态结果误差预算及阶段报告 & 动态参考、夹具、同步与算法误差分离 & P0\\
\bottomrule
\end{longtable}
\normalsize
\section{逐项方案编写接口（后续版本）}
本阶段\textbf{测试编号已预留，详细实施方案尚未完成}。每个编号在正式执行前，应按 S1 章节的模板补齐：目的、仪器和参考不确定度、进入条件、工况矩阵、台架操作、原始数据字段、计算公式、合格判据、异常处理和结果归档。不得因本表列有测试名称就自动判定已经执行或通过。
\begin{warningbox}正弦摆动 $\theta=A\sin(2\pi ft)$ 的峰值速率为 $2\pi fA$（$A$ 用度时速率是 $^\circ$/s）。需同时审核角加速度、转台摆动路径、轴系联锁、夹具偏心和安全限位。\end{warningbox}
